\documentclass[10pt,a4paper]{amsart}
\usepackage[margin=19mm]{geometry}
\usepackage{amsmath,amssymb,mathtools}
\usepackage[scr=rsfs,cal=euler]{mathalfa}
\usepackage{xfrac}
\usepackage[unicode,hidelinks,bookmarks=false]{hyperref}
\newcommand{\ie}{i.e.\ }
\newcommand{\cf}{cf.\ }
\newcommand{\resp}{respectively\ }
\newcommand{\Subsection}{\subsection}
\providecommand{\nobreakdash}{\nobreak\hbox{-}}
\setlength{\emergencystretch}{2em}
\multlinegap0pt
\usepackage{amssymb}
\usepackage[shortlabels]{enumitem}
\usepackage{esint}
\usepackage{tikz}
\newtheorem{proposition}{Proposition}
\newtheorem{lemma}{Lemma}
\newcommand{\cpct}{\operatorname{Cap}}
\newcommand{\csubset}{\subset\joinrel\subset}
\title{\textbf{Combined effect of homogenization and dimension-reduction in the random Neumann sieve problem}}
\author{\textbf{Mert Baştuğ}}
\date{Source: arXiv:2604.15183v1 (CC BY 4.0)}
\begin{document}
\maketitle

\noindent\emph{Source and licence.} \textbf{Combined effect of homogenization and dimension-reduction in the random Neumann sieve problem}. Version arXiv:2604.15183v1 (CC BY 4.0), distributed by arXiv under the Creative Commons Attribution 4.0 International licence, http://creativecommons.org/licenses/by/4.0/, as recorded in the arXiv licence metadata for this exact version and shown on the abstract page https://arxiv.org/abs/2604.15183v1. Full licence notice: \texttt{licenses/arxiv-2604.15183-CC-BY-4.0.txt}.\par\smallskip

\noindent\textit{Editorial scope.} Counted unit: the article's lemma lm:bad\_point\_big\_average (source0.tex lines 1448--1454). The article's complete original proof of it is delivered with it (source0.tex lines 1456--1539, one \texttt{proof} environment of 84 lines). No mathematics is written, repaired or re-derived: the statement and the proof are the article's own lines, in the article's own notation, and no retained line is substituted or re-flowed.  Retained uncounted as the source-local context the proof needs: lines 325--334; lines 421--426; lines 495--508; lines 510--519.  Nothing was omitted from any retained span, so the package carries no omission record.  No retained line contains a citation, so the wrapper carries no bibliography and no bibliography entry.  No retained line contains a figure environment, an image-inclusion command or drawing code, so no image is delivered and the package declares no asset dependency.  Editorial formatting only, in addition: an AMS \texttt{amsart} preamble that restates the article's own declarations and macros used by the retained lines, namely the environments \texttt{proposition}, \texttt{lemma} on separate counters, together with the standard packages those lines need.  The retained environments print in this document as Proposition 1, Lemma 1 in order of appearance, whereas the article prints the same items with the numbering of its own full document.  The complete native source of the article as distributed in the arXiv e-print for this exact version is bundled unchanged as \texttt{sources/198586/source0.tex}.
\begin{proposition} \label{pr:continuity_in_h}
    Let $K' \subset \mathbb{R}^{N - 1}$ be bounded, and let $U' \subset \mathbb{R}^{N - 1}$ be open with $\overline{K'} \subset U'$. If $h_1 \le h_2$, then
    \begin{equation} \label{eq:comparison_in_h}
        \frac{h_1}{h_2} \cpct_{h_2}(K', U') \le \cpct_{h_1}(K', U') \le \frac{h_2}{h_1} \cpct_{h_2}(K', U').
    \end{equation}
    In particular, $h \in (0, \infty) \mapsto \cpct_h(K', U')/(2h)$ is a continuous decreasing function. Furthermore,
    \begin{equation} \label{eq:finite_limit}
        \lim_{h \to 0} \frac{1}{2h} \cpct_h(K', U') < \infty.
    \end{equation}
\end{proposition}

\begin{proposition} \label{pr:limit_at_infinity}
    Assume $N > 3$. Let $K' \subset \mathbb{R}^{N - 1}$ be bounded, and let $U' \subset \mathbb{R}^{N - 1}$ be open with $\overline{K'} \subset U'$. Then $\cpct_h(K', U') \le \cpct(K' \times \{0\}, U' \times \mathbb{R})$ for all $h > 0$ and
    \[
    \lim_{h \to \infty} \cpct_h(K', U') = \cpct(K' \times \{0\}, U' \times \mathbb{R}).
    \]
\end{proposition}

\begin{proposition} \label{pr:subcapacitary_growth_1}
    Let $K \subset \mathbb{R}^3$ be bounded. Fix $\theta \in (0, 1)$. There exist positive constants $L$, $H$, $\tau$ and $c$ such that if
    \[
    l \ge L, \quad h \ge H, \quad \max \left\{\frac{\ln(l)}{h}, \frac{h}{l^2} \right\} \le \tau,
    \]
    and $v \in H^1(C(l, h))$ satisfies
    \[
    v = 1 \text{ weakly on } K, \quad \int_{C(l, h)} |\nabla v|^2 \, dx < \theta \cpct(K, \mathbb{R}^3),
    \]
    then we have the lower bound
    \[
    \fint_{C(l, h)} v \, dx > c.
    \]
\end{proposition}

\begin{proposition} \label{pr:subcapacitary_growth_2}
    Let $N > 3$ and let $K' \subset \mathbb{R}^{N - 1}$ be bounded. Fix $\theta \in (0, 1)$. There exist positive constants $L$, $\tau$ and $c$ such that if $l \ge L$, $h/l^{N - 1} \le \tau$, and $v \in H^1(C(l, h))$ satisfies
    \[
    v = 1 \text{ weakly on } K' \times \{0\}, \quad \int_{C(l, h)} |\nabla v|^2 \, dx < \theta \cpct_h(K', \mathbb{R}^{N - 1}),
    \]
    then we have the lower bound
    \[
    \fint_{C(l, h)} v \, dx > c.
    \]
\end{proposition}

\begin{lemma} \label{lm:bad_point_big_average}
    Let $\omega \in \Omega$ and $(v_\varepsilon) \in \mathcal{O}(\omega)$ be fixed. Let $\sigma, \theta \in (0, 1)$ and let $O' \csubset U'$ be open. There exist positive constants $c$ and $\varepsilon_0$ such that if $\varepsilon < \varepsilon_0$, then we have
    \begin{equation} \label{eq:bad_point_big_average}
        \fint_{B'(\varepsilon y', \varepsilon r(\omega, y')) \times (0, 1)} 1 - \hat{v}_\varepsilon \, dx > c
    \end{equation}
    for all $(y', \rho) \in \mathcal{N}_{\varepsilon, h_0}(\sigma, \theta, O')$.
\end{lemma}

\begin{proof}
    We extend $v_\varepsilon$ to $U' \times (-\delta_\varepsilon, \delta_\varepsilon)$ by reflection, that is, we set $v_\varepsilon(x', x_N) := v_\varepsilon(x', -x_N)$ for $(x', x_N) \in U' \times (-\delta_\varepsilon, 0)$. Then $v_\varepsilon \in H^1(U' \times (-\delta_\varepsilon, \delta_\varepsilon))$. If $(y', \rho) \in \mathcal{N}_{\varepsilon, h_0}(\sigma, \theta, O')$, we have
    \begin{equation} \label{eq:blow_up}
        \int_{B'(\varepsilon y', \varepsilon r(\omega, y')) \times (-\delta_\varepsilon, \delta_\varepsilon)} |\nabla v_\varepsilon|^2  \, dx < \theta \varepsilon^{N - 1} \delta_\varepsilon J_{h_0}(\rho).
    \end{equation}
    An easy computation shows that
    \begin{equation} \label{eq:capacity_identities}
    \varepsilon^{N - 1} \delta_\varepsilon J_{h_0}(\rho) = 
        \begin{cases}
            a_\varepsilon^{N - 2} \rho^{N - 2} \cpct(T' \times \{0\}, \mathbb{R}^N) \quad &\text{if } h_0 = \infty, \\
            a_\varepsilon^{N - 2} \rho^{N - 2} \cpct_{\frac{h_0}{\rho}}(T', \mathbb{R}^{N - 1}) \quad &\text{if } h_0 \in (0, \infty), \\
            2 a_\varepsilon^{N - 3} \delta_\varepsilon \rho^{N - 3} \cpct_0(T', \mathbb{R}^{N - 1}) \quad &\text{if } h_0 = 0.
        \end{cases}
    \end{equation}
    Let us now assume $N = 3$. In this case $h_0 = \infty$. Then \eqref{eq:blow_up} and \eqref{eq:capacity_identities} give
    \[
    \frac{1}{a_\varepsilon^{N - 2} \rho^{N - 2}} \int_{B'(\varepsilon y', \varepsilon r(\omega, y')) \times (-\delta_\varepsilon, \delta_\varepsilon)} |\nabla v_\varepsilon|^2 \, dx < \theta \cpct(T' \times \{0\}, \mathbb{R}^N).
    \]
    By applying a change of variables in Proposition \ref{pr:subcapacitary_growth_1}, we obtain positive constants $L$, $H$, $\tau$ and $c$, depending on $\theta$, such that if
    \begin{equation} \label{eq:desired_implication}
         \frac{\varepsilon r(\omega, y')}{a_\varepsilon \rho} \ge L, \quad \frac{h_\varepsilon}{\rho} \ge H, \quad \max \left\{\ln\left(\frac{\varepsilon r(\omega, y')}{a_\varepsilon \rho}\right) \frac{a_\varepsilon \rho}{\delta_\varepsilon}, \frac{h_\varepsilon}{\rho}\frac{a_\varepsilon^2 \rho^2}{\varepsilon^2 r(\omega, y')^2}\right\} \le \tau,
    \end{equation}
    then
    \begin{equation}
        c < \fint_{B'(\varepsilon y', \varepsilon r(\omega, y')) \times (-\delta_\varepsilon, \delta_\varepsilon)} 1 - v_\varepsilon \, dx = \fint_{B'(\varepsilon y', \varepsilon r(\omega, y')) \times (0, 1)} 1 - \hat{v}_\varepsilon \, dx.
    \end{equation}
    We show that there exists $\varepsilon_0$ such that the conditions in \eqref{eq:desired_implication} are satisfied for all $(y', \rho) \in \mathcal{N}_{\varepsilon, h_0}(\sigma, \theta, O')$ if $\varepsilon < \varepsilon_0$. We write the first two conditions as
    \begin{equation} \label{eq:equivalent_bounds}
        \frac{\varepsilon}{a_\varepsilon} \ge \frac{L \rho}{r(\omega, y')}, \quad h_\varepsilon \ge H \rho.
    \end{equation}
    Since $(y', \rho) \in M(\omega) \setminus M_\sigma(\omega)$, we know that $r(\omega, y') \ge \sigma$ and $\rho \le 1/\sigma$. Hence, the following conditions imply \eqref{eq:equivalent_bounds} and are independent of $(y', \rho)$:
    \begin{equation} \label{eq:equivalent_bounds_independent}
        \frac{\varepsilon}{a_\varepsilon} \ge \frac{L}{\sigma^2}, \quad h_\varepsilon \ge \frac{H}{\sigma}.
    \end{equation}
    Next, we rewrite the third condition in \eqref{eq:desired_implication} as
    \begin{equation} \label{eq:equivalent_bounds_2}
        \left(\ln\left(\frac{r(\omega, y')}{\varepsilon \rho}\right) + \ln\left(\frac{1}{\delta_\varepsilon}\right) \right) \varepsilon^2 \rho \le \tau, \quad \frac{\rho}{r(\omega, y')} \delta_\varepsilon^2 \le \tau.
    \end{equation}
    We know that $x \ln(1/x)$ is increasing for $x \le 1/e$. Hence, using $r(\omega, y') \le 1$ and $\rho \le 1/\sigma$, we get
    \[
    \varepsilon^2 \rho \ln\left(\frac{r(\omega, y')}{\varepsilon \rho}\right) + \varepsilon^2 \rho \ln\left(\frac{1}{\delta_\varepsilon}\right) \le \frac{1}{\sigma} \left(\varepsilon^2 \ln\left(\frac{\sigma}{\varepsilon}\right) + \varepsilon^2 \ln\left(\frac{1}{\delta_\varepsilon}\right)\right)
    \]
    for $\varepsilon \le \sigma/e$. Thus
    \begin{equation} \label{eq:equivalent_bounds_independent_2}
        \tau \ge \frac{1}{\sigma} \left(\varepsilon^2 \ln\left(\frac{\sigma}{\varepsilon}\right) + \varepsilon^2 \ln\left(\frac{1}{\delta_\varepsilon}\right)\right)
    \end{equation}
    implies the first inequality in \eqref{eq:equivalent_bounds_2} for sufficiently small $\varepsilon$ and is independent of $(y', \rho)$. On the other hand,
    \begin{equation} \label{eq:equivalent_bounds_independent_3}
        \tau \ge \frac{\delta_\varepsilon^2}{\sigma^3}
    \end{equation}
    implies the second inequality in \eqref{eq:equivalent_bounds_2} and is independent of $(y', \rho)$. To conclude the proof for $N = 3$, we need to show that \eqref{eq:equivalent_bounds_independent}, \eqref{eq:equivalent_bounds_independent_2} and \eqref{eq:equivalent_bounds_independent_3} are all satisfied for sufficiently small $\varepsilon$. However, this follows easily from the following limits:
    \[
    \lim_{\varepsilon \to 0} \varepsilon^2 \ln\left(\frac{1}{\delta_\varepsilon}\right) = \lim_{\varepsilon \to 0} \delta_\varepsilon = 0, \quad \lim_{\varepsilon \to 0} \frac{\varepsilon}{a_\varepsilon} = \lim_{\varepsilon \to 0} h_\varepsilon = \infty.
    \]
    The first limit holds by assumption, and the last limit holds because $h_0 = \infty$.

    Let $N > 3$. We prove that for any $\tilde{\theta} \in (\theta, 1)$ and for sufficiently small $\varepsilon$, we have
    \begin{equation} \label{eq:change_of_capacity_function}
        \frac{1}{a_\varepsilon^{N - 2} \rho^{N - 2}} \int_{B'(\varepsilon y', \varepsilon r(\omega, y')) \times (-\delta_\varepsilon, \delta_\varepsilon)} |\nabla v_\varepsilon|^2 \, dx \le \tilde{\theta} \cpct_{\frac{h_\varepsilon}{\rho}}(T', \mathbb{R}^{N - 1})
    \end{equation}
    for all $(y', \rho) \in \mathcal{N}_{\varepsilon, h_0}(\sigma, \theta, O')$. Given the claim, we can use Proposition \ref{pr:subcapacitary_growth_2} to obtain $L$, $c$ and $\tau$, depending on $\tilde{\theta}$, such that if
    \begin{equation} \label{eq:desired_implication_2}
        \frac{\varepsilon r(\omega, y')}{a_\varepsilon \rho} \ge L, \quad \frac{h_\varepsilon}{\rho} \left(\frac{a_\varepsilon \rho}{\varepsilon r(\omega, y')}\right)^{N - 1} \le \tau,
    \end{equation}
    then
    \[
    c < \fint_{B'(\varepsilon y', \varepsilon r(\omega, y')) \times (-\delta_\varepsilon, \delta_\varepsilon)} 1 - v_\varepsilon \, dx = \fint_{B'(\varepsilon y', \varepsilon r(\omega, y')) \times (0, 1)} 1 - \hat{v}_\varepsilon \, dx.
    \]
    We can establish, as in the case $N = 3$, that there exists $\varepsilon_0$ such that if $\varepsilon < \varepsilon_0$, then the conditions of \eqref{eq:desired_implication_2} are satisfied for all $(y', \rho) \in \mathcal{N}_{\varepsilon, h_0}(\sigma, \theta, O')$. We do not carry out the computations here and turn to the proof of \eqref{eq:change_of_capacity_function} instead.

    Let $\tilde{\theta} \in (\theta, 1)$  be arbitrary. Assume $h_0 = \infty$. By Proposition \ref{pr:limit_at_infinity}, we have that $\theta \cpct(T' \times \{0\}, \mathbb{R}^N) \le \tilde{\theta} \cpct_h(T', \mathbb{R}^{N - 1})$ for sufficiently large $h$. Hence, \eqref{eq:blow_up} and \eqref{eq:capacity_identities} give
    \begin{equation*}
        \frac{1}{a_\varepsilon^{N - 2} \rho^{N - 2}} \int_{B'(\varepsilon y', \varepsilon r(\omega, y')) \times (-\delta_\varepsilon, \delta_\varepsilon)} |\nabla v_\varepsilon|^2 \, dx \le \theta \cpct(T' \times \{0\}, \mathbb{R}^N) \le \tilde{\theta} \cpct_{\frac{h_\varepsilon}{\rho}}(T', \mathbb{R}^{N - 1}).
    \end{equation*}
    for small enough $\varepsilon$. Next, assume $h_0 = 0$. Since $\cpct_h(T', \mathbb{R}^{N - 1})/(2h) \to \cpct_0(T', \mathbb{R}^{N - 1})$ as $h \to 0$ by definition, we know that $\theta \, 2h \cpct_0(T', \mathbb{R}^{N - 1}) \le \tilde{\theta} \cpct_h(T', \mathbb{R}^{N - 1})$ for sufficiently small $h$. Therefore, by \eqref{eq:blow_up} and \eqref{eq:capacity_identities},
    \begin{equation*}
        \frac{1}{a_\varepsilon^{N - 2} \rho^{N - 2}} \int_{B'(\varepsilon y', \varepsilon r(\omega, y')) \times (-\delta_\varepsilon, \delta_\varepsilon)} |\nabla v_\varepsilon|^2 \, dx \le \theta \frac{2h_\varepsilon}{\rho} \cpct_0(T', \mathbb{R}^{N - 1}) \le \tilde{\theta} \cpct_{\frac{h_\varepsilon}{\rho}}(T', \mathbb{R}^{N - 1})
    \end{equation*}
    for small enough $\varepsilon$. Finally, let $h_0 \in (0, \infty)$. Pick $\varepsilon$ small enough such that $\theta \max \{h_\varepsilon/h_0, h_0/h_\varepsilon\} \le \tilde{\theta}$. Then, Proposition \ref{pr:continuity_in_h}, together with \eqref{eq:blow_up} and \eqref{eq:capacity_identities}, implies
    \begin{equation*}
        \frac{1}{a_\varepsilon^{N - 2} \rho^{N - 2}} \int_{B'(\varepsilon y', \varepsilon r(\omega, y')) \times (-\delta_\varepsilon, \delta_\varepsilon)} |\nabla v_\varepsilon|^2 \, dx < \theta \cpct_{\frac{h_0}{\rho}}(T', \mathbb{R}^{N - 1}) \le \tilde{\theta} \cpct_\frac{h_\varepsilon}{\rho}(T', \mathbb{R}^{N - 1}).
    \end{equation*}
    This concludes the proof of the inequality \eqref{eq:change_of_capacity_function}.
\end{proof}

\end{document}
